\section{Sensitivity to Temperature and Number of Samples}
\label{app:sensitivity}

In \Cref{sec:SE_AST}, we evaluate multi-sample UQ methods using the default settings of \citet{farquhar2024detec}: temperature $T\!=\!1.0$ and $J\!=\!10$ samples. Here, we ablate these choices for the best-performing multi-sample method, SE$_{AST}$.

\Cref{tab:sensitivity_temperature} shows the change in mean AUROC (over all models) when varying $T$ while keeping $J\!=\!10$ fixed. Lowering $T$ to $0.5$ has a detrimental effect (up to $-0.09$), while increasing to $T\!=\!1.5$ can improve AUROC by up to $0.04$. However, increasing $T$ further undoes these improvements and can even lead to a decrease in AUROC compared to $T\!=\!1.0$. For splits involving the \emph{Irrelevance} task, increasing $T$ to $1.5$ yields almost no improvement and becomes detrimental at higher temperatures as well.

\Cref{tab:sensitivity_nsamples} shows the change in mean AUROC when varying $J$ while keeping $T\!=\!1.0$ fixed. Increasing $J$ to $30$ can improve AUROC by around $0.02$--$0.03$, but improvements flatten out beyond that (no better performance for $J\!=\!40$).

Importantly, while the improvements from tuning $T$ and $J$ can close the gap to G-NLL and G-NLL$_{SMT}$, except for the \emph{Multiple} split they do not surpass them. In light of the added computational cost required for multiple samples (\Cref{sec:SE_AST}), our main conclusion remains unchanged: SE as the best multi-sample UQ method offers no clear advantage over G-NLL as the best single-sample method.


\begin{table}[h!]
\centering
\caption{Differences of mean AUROC of SE$_{AST}$ w.r.t.\ its values for $T\!=\!1.0$. $J\!=\!10$ fixed for all.}
\label{tab:sensitivity_temperature}
\begin{tabular}{l|ccccc}
\toprule
\textbf{Task} & $T\!=\!0.5$ & $T\!=\!1.0$ & $T\!=\!1.5$ & $T\!=\!2.0$ & $T\!=\!3.0$ \\
\midrule
Simple & $-0.09$ & $0$ & $0.03$ & $0.03$ & $0.02$ \\
Multiple & $-0.06$ & $0$ & $0.03$ & $0.02$ & $0$ \\
Parallel & $-0.01$ & $0$ & $0.04$ & $-0.01$ & $0.02$ \\
Parallel-Multiple & $-0.05$ & $0$ & $0.03$ & $-0.03$ & $-0.03$ \\
Simple + Multiple & $-0.07$ & $0$ & $0.03$ & $0.03$ & $0.01$ \\
Simple + Parallel & $-0.05$ & $0$ & $0.04$ & $0.01$ & $0.01$ \\
Multiple + Parallel-Multiple & $-0.06$ & $0$ & $0.03$ & $-0.01$ & $-0.03$ \\
All Combined & $-0.05$ & $0$ & $0.04$ & $0$ & $-0.01$ \\
Simple + Irrelevance & $-0.07$ & $0$ & $0$ & $0$ & $-0.03$ \\
All Combined + Irrelevance & $-0.06$ & $0$ & $0.01$ & $-0.01$ & $-0.03$ \\
\bottomrule
\end{tabular}
\end{table}


\begin{table}[h!]
\centering
\caption{Differences of mean AUROC of SE$_{AST}$ w.r.t.\ its values for $J\!=\!10$. $T\!=\!1.0$ fixed for all.}
\label{tab:sensitivity_nsamples}
\begin{tabular}{l|ccccc}
\toprule
\textbf{Task} & $J\!=\!5$ & $J\!=\!10$ & $J\!=\!20$ & $J\!=\!30$ & $J\!=\!40$ \\
\midrule
Simple & $-0.03$ & $0$ & $0.01$ & $0.02$ & $0.02$ \\
Multiple & $-0.05$ & $0$ & $0.01$ & $0.01$ & $0.01$ \\
Parallel & $-0.04$ & $0$ & $0.02$ & $0.02$ & $0.02$ \\
Parallel-Multiple & $-0.03$ & $0$ & $0.01$ & $0.03$ & $0.03$ \\
Simple + Multiple & $-0.04$ & $0$ & $0.01$ & $0.02$ & $0.02$ \\
Simple + Parallel & $-0.03$ & $0$ & $0.02$ & $0.02$ & $0.02$ \\
Multiple + Parallel-Multiple & $-0.04$ & $0$ & $0.01$ & $0.03$ & $0.03$ \\
All Combined & $-0.03$ & $0$ & $0.02$ & $0.03$ & $0.03$ \\
Simple + Irrelevance & $-0.02$ & $0$ & $0.01$ & $0.02$ & $0.02$ \\
All Combined + Irrelevance & $-0.03$ & $0$ & $0.01$ & $0.02$ & $0.02$ \\
\bottomrule
\end{tabular}
\end{table}


\section{Paired Significance Analysis}
\label{app:paired_significance}
We perform a paired statistical analysis for G-NLL$_{SMT}$ vs G-NLL, and SE$_{AST}$ vs SE$_{EXM}$ to test whether our proposed SMT and AST clustering improve over their baselines. Additionally, we also include a test between the best single-sample method and the best multi-sample method, G-NLL vs SE$_{AST}$. In each pair the first-named method is method A, so that a positive $\Delta$ favours the proposed method. For each task and each pair of UQ methods, we test whether one method's AUROC is systematically higher across models. For every one of the eight LLMs we compute each method's AUROC, giving one paired difference per model, $\Delta$AUROC = AUROC(A) - AUROC(B). We apply a two-sided Wilcoxon signed-rank test to these eight paired differences, separately for each task. We report the median across model $\Delta$AUROC as the effect size together with the two-sided p-value (* denotes p < 0.05). A positive $\Delta$ favors method A. Note that with eight models, the statistical power is limited, with the smallest attainable two-sided $p = \frac{2}{2^{8}} \approx 0.008$, which is reached only when all eight per-model differences share the same sign. Therefore, the effect sizes (median $\Delta$AUROC) should be considered alongside $p$-values when interpreting the results.

As shown in Table~\ref{tab:wilcoxon}, G-NLL$_{SMT}$ vs G-NLL shows mixed results: the median $\Delta$AUROC is small on every task ($|\Delta| \leq 0.017$) and split in both directions, being negative on four tasks, positive on four, and vanishing to three decimal places on the remaining two, with no task reaching significance. We therefore find no statistical evidence that SMT improves over G-NLL across models. In contrast, SE$_{AST}$ vs.\ SE$_{EXM}$ is positive on nine of the ten task splits, indicating that AST-based clustering consistently improves over EXM-based clustering; the sole exception is \emph{Simple + Irrelevance}, where AST-based clustering is marginally worse ($\Delta = -0.004$, $p = 0.383$). The effect is modest in size but consistent in direction, reaching significance on \emph{Parallel} ($\Delta = +0.014$, $p = 0.023$) and falling just short of it on the largest split (\emph{All Combined}, $\Delta = +0.020$, $p = 0.055$). Finally, comparing the best single-sample method (G-NLL) against the best multi-sample method (SE$_{AST}$), the median $\Delta$AUROC favours G-NLL on all ten splits, significantly so on both unanswerable splits (\emph{Simple + Irrelevance}, $\Delta = +0.061$, $p = 0.008$; \emph{All Combined + Irrelevance}, $\Delta = +0.040$, $p = 0.039$); on the former, $p$ attains the minimum possible value for eight models, i.e.\ G-NLL is better on every model. This shows that G-NLL, a single-sample greedy method, matches or exceeds the strongest sampling-based estimator, at a fraction of its inference cost.



\begin{table}[h!]
\centering
\scalebox{0.62}{
\begin{tabular}{llccc}
\toprule
\textbf{Category} & \textbf{Task} & \textbf{G-NLL$_{\mathrm{SMT}}$ vs G-NLL} & \textbf{SE$_{\mathrm{AST}}$ vs SE$_{\mathrm{EXM}}$} & \textbf{G-NLL vs SE$_{\mathrm{AST}}$} \\
\midrule
\multirow{4}{*}{Individual} & Simple & $-0.012\phantom{^{*}}$ (0.844) & $+0.022\phantom{^{*}}$ (0.312) & $+0.023\phantom{^{*}}$ (0.383) \\
 & Multiple & $-0.004\phantom{^{*}}$ (0.945) & $+0.029\phantom{^{*}}$ (0.250) & $+0.006\phantom{^{*}}$ (0.641) \\
 & Parallel & $+0.000\phantom{^{*}}$ (0.578) & $+0.014^{*}$ (0.023) & $+0.033\phantom{^{*}}$ (0.312) \\
 & Parallel-Multiple & $+0.017\phantom{^{*}}$ (0.312) & $+0.020\phantom{^{*}}$ (0.461) & $+0.032\phantom{^{*}}$ (0.461) \\
\midrule
\multirow{4}{*}{Combinations} & Simple + Multiple & $-0.005\phantom{^{*}}$ (0.742) & $+0.020\phantom{^{*}}$ (0.148) & $+0.018\phantom{^{*}}$ (0.461) \\
 & Simple + Parallel & $-0.012\phantom{^{*}}$ (0.844) & $+0.023\phantom{^{*}}$ (0.109) & $+0.018\phantom{^{*}}$ (0.383) \\
 & Multiple + Parallel-Multiple & $+0.009\phantom{^{*}}$ (0.312) & $+0.024\phantom{^{*}}$ (0.195) & $+0.019\phantom{^{*}}$ (0.461) \\
 & All Combined & $-0.000\phantom{^{*}}$ (0.641) & $+0.020\phantom{^{*}}$ (0.055) & $+0.021\phantom{^{*}}$ (0.312) \\
\midrule
\multirow{2}{*}{Unanswerable} & Simple + Irrelevance & $+0.007\phantom{^{*}}$ (0.250) & $-0.004\phantom{^{*}}$ (0.383) & $+0.061^{*}$ (0.008) \\
 & All Combined + Irrelevance & $+0.005\phantom{^{*}}$ (0.383) & $+0.005\phantom{^{*}}$ (0.312) & $+0.040^{*}$ (0.039) \\
\bottomrule
\end{tabular}
}
\caption{Paired two-sided Wilcoxon signed-rank tests across the eight LLMs. Each cell reports the median across-model $\Delta$AUROC (first method minus second) and its $p$-value; $^{*}$ marks $p<0.05$. A positive $\Delta$ favors the first-named method. Columns: (i) whether SMT improves G-NLL, (ii) whether AST clustering improves semantic entropy, and (iii) the best single-sample (G-NLL) vs.\ the best multi-sample (SE$_{\mathrm{AST}}$) method.}
\label{tab:wilcoxon}

\end{table}



\section{Black-box Methods Comparison}
\label{app:black_box}
Beyond the DSE-based methods, which require no log-probabilities, we add self-consistency (SC), or the majority-cluster frequency across the multiple samples under both clustering schemes, in Table~\ref{tab:blackbox}. All four scores are computed from ten sampled generations and their semantic clustering only, without access to token log-probabilities. A higher SC score indicates lower uncertainty. We note that the difference between SC and DSE is tiny, while SC slightly underperforms DSE. This suggests that the black-box signal comes mostly from the clustering, rather than from the choice between an entropy and a majority-vote functional.

\begin{table}[h!]
\centering
\scalebox{0.7}{
\begin{tabular}{llcccc}
\toprule
& & \multicolumn{2}{c}{\textbf{EXM clustering}} & \multicolumn{2}{c}{\textbf{AST clustering}} \\
\cmidrule(lr){3-4} \cmidrule(lr){5-6}
\textbf{Category} & \textbf{Task} & $\mathbf{DSE}_{EXM}$ & $\mathbf{SC}_{EXM}$ & $\mathbf{DSE}_{AST}$ & $\mathbf{SC}_{AST}$ \\
\midrule
\multirow{4}{*}{Individual} & Simple & $0.73$ \tiny{±$.050$} & $0.73$ \tiny{±$.049$} & \best{$0.75$} \tiny{±$.047$} & \underline{$0.75$} \tiny{±$.046$} \\
 & Multiple & $0.73$ \tiny{±$.073$} & $0.72$ \tiny{±$.073$} & \best{$0.74$} \tiny{±$.067$} & \underline{$0.74$} \tiny{±$.066$} \\
 & Parallel & $0.66$ \tiny{±$.059$} & $0.66$ \tiny{±$.058$} & \best{$0.69$} \tiny{±$.056$} & \underline{$0.68$} \tiny{±$.055$} \\
 & Parallel-Multiple & $0.73$ \tiny{±$.051$} & $0.72$ \tiny{±$.052$} & \best{$0.73$} \tiny{±$.050$} & \underline{$0.73$} \tiny{±$.049$} \\
\midrule
\multirow{4}{*}{Combinations} & Simple + Multiple & $0.73$ \tiny{±$.041$} & $0.72$ \tiny{±$.040$} & \best{$0.75$} \tiny{±$.038$} & \underline{$0.74$} \tiny{±$.038$} \\
 & Simple + Parallel & $0.70$ \tiny{±$.038$} & $0.70$ \tiny{±$.037$} & \best{$0.73$} \tiny{±$.036$} & \underline{$0.72$} \tiny{±$.036$} \\
 & Multiple + Parallel-Multiple & $0.73$ \tiny{±$.041$} & $0.73$ \tiny{±$.040$} & \best{$0.75$} \tiny{±$.038$} & \underline{$0.75$} \tiny{±$.039$} \\
 & All Combined & $0.72$ \tiny{±$.028$} & $0.72$ \tiny{±$.027$} & \best{$0.74$} \tiny{±$.026$} & \underline{$0.74$} \tiny{±$.026$} \\
\midrule
\multirow{2}{*}{Unanswerable} & Simple + Irrelevance & \best{$0.76$} \tiny{±$.026$} & \best{$0.76$} \tiny{±$.026$} & $0.74$ \tiny{±$.026$} & $0.73$ \tiny{±$.026$} \\
 & All Combined + Irrelevance & \best{$0.74$} \tiny{±$.021$} & $0.73$ \tiny{±$.020$} & \underline{$0.73$} \tiny{±$.020$} & $0.73$ \tiny{±$.020$} \\
\bottomrule
\end{tabular}
}
\caption{Black-box uncertainty methods: Mean AUROC $\pm$ mean standard error (mean over eight LLMs). All four scores are computed from ten sampled generations and their semantic clustering only, without access to token log-probabilities: discrete semantic entropy (DSE) and self-consistency (SC; the frequency of the majority cluster), each under exact-match (EXM) and AST-based clustering. The best result in \best{green}, \underline{underlined} second best. Note the slight variation of numbers compared to \Cref{tab:combined_all_tasks} and \Cref{tab:new_method_se} are due to high temperature sampling variations.}
\label{tab:blackbox}
\end{table}

\section{Selective Prediction Metrics}
\label{app:selection_pred_metrics}
Metrics such as AUROC as shown in Table~\ref{tab:combined_all_tasks} measure discrimination, or how well an uncertainty score ranks correct over incorrect predictions, but they do not reflect the utility of the score under abstention. Here we report two selective prediction metrics, the normalized area under the risk-coverage curve (NAURC) and the prediction rejection ratio (PRR).

\paragraph{NAURC.} The area under the risk-coverage curve (AURC) summarizes selective prediction performance across all abstention thresholds from a coverage $c$ of 0 to 1. The model answers the $c$ fraction of inputs with the lowest uncertainty and we compute the corresponding selective risk, which is the error rate among the answered inputs. AURC integrates this risk over coverage. As AURC depends on the model's base error rate, we report a normalized version of this metric, NAURC~\citep{cattelan2024fixbroken}:
$$\mathrm{NAURC} = \frac{\mathrm{AURC} - \mathrm{AURC}^{\star}}{\mathrm{AURC}_{\text{rand}} - \mathrm{AURC}^{\star}},$$
where $\text{AURC}^{\star}$ is the oracle AURC and $\text{AURC}_{\text{rand}}$ is the risk under random ordering (equal to the base error rate). A lower NAURC is better.

\paragraph{PRR.} We also report the prediction-rejection ratio (PRR). For a given rejection rate $k$, PRR$@k$ measures how much of the risk reduction achieved by an oracle ranking an uncertainty score can realize. i.e.\ the fraction of the maximal achievable error reduction that the score captures when rejecting up to the $k\%$ most uncertain input:
$$  \mathrm{PRR}@k =
  \frac{A_{\mathrm{random}} - A_{\mathrm{uncertainty}}}
       {A_{\mathrm{random}} - A_{\mathrm{oracle}}}.$$
A PRR of 1 matches the oracle, while 0 matches random rejection. To give an example, a PRR$@k$ of 0.5 means of all the error reduction that was achievable by rejecting up to $k$ percent of examples, the uncertainty score captures 50 percent. Here we report PRR$@30$ and PRR$@50$, which integrate the error-rejection curve as rejection increases from $0$ up to $30\%$ and $50\%$.


We report NAURC in Table~\ref{tab:naurc} and PRR in Tables~\ref{tab:prr30} and~\ref{tab:prr50}. For NAURC, we find a similar trend to the AUROC result in Table~\ref{tab:combined_all_tasks}: the single-sample scores dominate, with G-NLL$_{SMT}$ the best method on 9 of the 10 splits and G-NLL the best of the remaining methods, while every multi-sample score is clearly behind. The PRR tables refine this picture in a way the AUROC results do not reveal. At PRR$@30\%$, $\text{SE}_{AST}$ outperforms plain G-NLL on all eight answerable splits, so between the two unadapted scores the multi-sample one is indeed preferable in the low-rejection regime; this is exactly the regime in which a deployed system would abstain. Once the FC-adapted single-sample score is included, however, G-NLL$_{SMT}$ recovers most of that gap: it is the best method on 7 of the 10 splits at PRR$@30\%$, with $\text{SE}_{AST}$ retaining an advantage only on \emph{Multiple}, \emph{Simple + Multiple} and \emph{Simple + Parallel}, and it is the best method on all 10 splits at PRR$@50\%$. The FC-specific adaptation therefore matters most precisely where the generic single-sample score was weakest.

Two observations are worth separating here. First, the ordering between G-NLL$_{SMT}$ and G-NLL is far more consistent under these selective-prediction metrics than under AUROC: SMT improves PRR on all ten splits at both rejection rates and NAURC on nine, whereas its AUROC differences are small and split in both directions (\Cref{app:paired_significance}). This is not a contradiction. AUROC summarizes the ranking over the whole score range, including regions no deployed system operates in, whereas NAURC and PRR$@k$ weight the low-uncertainty region where abstention decisions are actually made. A modification that sharpens the ordering among the most uncertain outputs, which is what restricting the score to semantically meaningful tokens is intended to do, can therefore move PRR appreciably while leaving AUROC almost unchanged. Second, we report effect sizes and the consistency of their direction rather than $p$-values for these two metrics. We deliberately do not convert the split-level win counts into a test: the ten splits are overlapping combinations of the same four base task sets, so they are not independent observations and a count-based $p$-value over splits would be invalid. This is precisely why the paired test in \Cref{app:paired_significance} is taken across models rather than across splits. We note finally that NAURC values greater than 1, as well as a negative PRR$@50\%$ for LEN, indicate performance worse than random.


\begin{table}[h!]
\centering
\scalebox{0.5}{
\begin{tabular}{llcccccccccc}
\toprule
& & \multicolumn{4}{c}{\textbf{Single-sample}} & \multicolumn{4}{c}{\textbf{Multi-sample}} & \multicolumn{2}{c}{\textbf{Other}} \\
\cmidrule(lr){3-6} \cmidrule(lr){7-10} \cmidrule(lr){11-12}
\textbf{Category} & \textbf{Task} & \textbf{MAX} & \textbf{AVG} & \textbf{G-NLL} & \textbf{G-NLL}$_{\mathrm{SMT}}$ & $\mathbf{SE}_{EXM}$ & $\mathbf{SE}_{AST}$ & $\mathbf{DSE}_{EXM}$ & \textbf{PE} & $\mathbf{P}_{true}$ & \textbf{LEN} \\
\midrule
\multirow{4}{*}{Individual} & Simple & $0.41$ \tiny{±$.098$} & $0.39$ \tiny{±$.103$} & \best{$0.36$} \tiny{±$.097$} & \best{$0.36$} \tiny{±$.113$} & $0.47$ \tiny{±$.145$} & $0.40$ \tiny{±$.150$} & $0.45$ \tiny{±$.144$} & $0.45$ \tiny{±$.106$} & $0.54$ \tiny{±$.128$} & $0.84$ \tiny{±$.155$} \\
 & Multiple & $0.46$ \tiny{±$.144$} & $0.42$ \tiny{±$.142$} & \underline{$0.39$} \tiny{±$.136$} & \best{$0.38$} \tiny{±$.129$} & $0.56$ \tiny{±$.209$} & $0.49$ \tiny{±$.211$} & $0.56$ \tiny{±$.212$} & $0.50$ \tiny{±$.160$} & $0.59$ \tiny{±$.182$} & $0.83$ \tiny{±$.238$} \\
 & Parallel & $0.59$ \tiny{±$.143$} & \underline{$0.48$} \tiny{±$.125$} & $0.49$ \tiny{±$.125$} & \best{$0.45$} \tiny{±$.121$} & $0.72$ \tiny{±$.180$} & $0.68$ \tiny{±$.177$} & $0.73$ \tiny{±$.181$} & $0.56$ \tiny{±$.136$} & $0.74$ \tiny{±$.163$} & $1.15$ \tiny{±$.215$} \\
 & Parallel-Multiple & $0.45$ \tiny{±$.101$} & $0.42$ \tiny{±$.104$} & \underline{$0.38$} \tiny{±$.096$} & \best{$0.34$} \tiny{±$.083$} & $0.49$ \tiny{±$.144$} & $0.47$ \tiny{±$.144$} & $0.49$ \tiny{±$.142$} & $0.45$ \tiny{±$.102$} & $0.57$ \tiny{±$.122$} & $0.87$ \tiny{±$.160$} \\
\midrule
\multirow{4}{*}{Combinations} & Simple + Multiple & $0.42$ \tiny{±$.080$} & $0.40$ \tiny{±$.081$} & \best{$0.36$} \tiny{±$.078$} & \underline{$0.37$} \tiny{±$.090$} & $0.51$ \tiny{±$.118$} & $0.51$ \tiny{±$.121$} & $0.53$ \tiny{±$.120$} & $0.47$ \tiny{±$.088$} & $0.56$ \tiny{±$.109$} & $0.85$ \tiny{±$.136$} \\
 & Simple + Parallel & $0.47$ \tiny{±$.084$} & $0.49$ \tiny{±$.088$} & \best{$0.39$} \tiny{±$.075$} & \best{$0.39$} \tiny{±$.080$} & $0.58$ \tiny{±$.111$} & $0.59$ \tiny{±$.116$} & $0.60$ \tiny{±$.113$} & $0.54$ \tiny{±$.091$} & $0.53$ \tiny{±$.087$} & $0.77$ \tiny{±$.111$} \\
 & Multiple + Parallel-Multiple & $0.43$ \tiny{±$.081$} & $0.47$ \tiny{±$.086$} & \underline{$0.36$} \tiny{±$.078$} & \best{$0.32$} \tiny{±$.069$} & $0.49$ \tiny{±$.120$} & $0.50$ \tiny{±$.118$} & $0.53$ \tiny{±$.120$} & $0.51$ \tiny{±$.087$} & $0.54$ \tiny{±$.101$} & $0.68$ \tiny{±$.127$} \\
 & All Combined & $0.45$ \tiny{±$.059$} & $0.48$ \tiny{±$.061$} & \underline{$0.37$} \tiny{±$.054$} & \best{$0.35$} \tiny{±$.056$} & $0.53$ \tiny{±$.081$} & $0.52$ \tiny{±$.081$} & $0.54$ \tiny{±$.082$} & $0.52$ \tiny{±$.064$} & $0.54$ \tiny{±$.070$} & $0.72$ \tiny{±$.086$} \\
\midrule
\multirow{2}{*}{Unanswerable} & Simple + Irrelevance & $0.37$ \tiny{±$.058$} & \underline{$0.33$} \tiny{±$.055$} & \underline{$0.33$} \tiny{±$.055$} & \best{$0.32$} \tiny{±$.060$} & $0.41$ \tiny{±$.075$} & $0.42$ \tiny{±$.080$} & $0.51$ \tiny{±$.076$} & $0.36$ \tiny{±$.058$} & $0.44$ \tiny{±$.066$} & $1.07$ \tiny{±$.106$} \\
 & All Combined + Irrelevance & $0.42$ \tiny{±$.045$} & $0.38$ \tiny{±$.044$} & \underline{$0.37$} \tiny{±$.042$} & \best{$0.34$} \tiny{±$.043$} & $0.48$ \tiny{±$.060$} & $0.51$ \tiny{±$.061$} & $0.51$ \tiny{±$.059$} & $0.42$ \tiny{±$.046$} & $0.52$ \tiny{±$.054$} & $1.09$ \tiny{±$.082$} \\
\bottomrule
\end{tabular}
}
\caption{Mean NAURC $\pm$ mean standard error (mean over eight LLMs) for individual tasks, task combinations, and unanswerable requests described in \Cref{sec:individual_tasks}, \Cref{sec:combinations}, and \Cref{sec:irrelevance_combinations} respectively. \textbf{Lower is better}; best result in \best{green}, \underline{underlined} second best.}
\label{tab:naurc}
\end{table}


\begin{table}[h!]
\centering
\scalebox{0.5}{
\begin{tabular}{llcccccccccc}
\toprule
& & \multicolumn{4}{c}{\textbf{Single-sample}} & \multicolumn{4}{c}{\textbf{Multi-sample}} & \multicolumn{2}{c}{\textbf{Other}} \\
\cmidrule(lr){3-6} \cmidrule(lr){7-10} \cmidrule(lr){11-12}
\textbf{Category} & \textbf{Task} & \textbf{MAX} & \textbf{AVG} & \textbf{G-NLL} & \textbf{G-NLL}$_{\mathrm{SMT}}$ & $\mathbf{SE}_{EXM}$ & $\mathbf{SE}_{AST}$ & $\mathbf{DSE}_{EXM}$ & \textbf{PE} & $\mathbf{P}_{true}$ & \textbf{LEN} \\
\midrule
\multirow{4}{*}{Individual} & Simple & $0.29$ \tiny{±$.091$} & $0.37$ \tiny{±$.092$} & $0.39$ \tiny{±$.094$} & \best{$0.43$} \tiny{±$.092$} & $0.38$ \tiny{±$.096$} & \underline{$0.42$} \tiny{±$.094$} & $0.38$ \tiny{±$.095$} & $0.29$ \tiny{±$.089$} & $0.28$ \tiny{±$.092$} & $0.05$ \tiny{±$.075$} \\
 & Multiple & $0.26$ \tiny{±$.122$} & $0.34$ \tiny{±$.130$} & $0.39$ \tiny{±$.134$} & \underline{$0.40$} \tiny{±$.129$} & $0.38$ \tiny{±$.133$} & \best{$0.43$} \tiny{±$.132$} & $0.38$ \tiny{±$.137$} & $0.31$ \tiny{±$.128$} & $0.24$ \tiny{±$.124$} & $0.13$ \tiny{±$.123$} \\
 & Parallel & $0.24$ \tiny{±$.107$} & $0.36$ \tiny{±$.109$} & $0.35$ \tiny{±$.111$} & \best{$0.38$} \tiny{±$.113$} & $0.32$ \tiny{±$.113$} & \underline{$0.37$} \tiny{±$.117$} & $0.31$ \tiny{±$.112$} & $0.28$ \tiny{±$.113$} & $0.17$ \tiny{±$.110$} & $0.00$ \tiny{±$.089$} \\
 & Parallel-Multiple & $0.31$ \tiny{±$.106$} & $0.37$ \tiny{±$.105$} & $0.39$ \tiny{±$.106$} & \best{$0.44$} \tiny{±$.105$} & $0.35$ \tiny{±$.111$} & \underline{$0.42$} \tiny{±$.102$} & $0.35$ \tiny{±$.111$} & $0.33$ \tiny{±$.114$} & $0.29$ \tiny{±$.121$} & $0.09$ \tiny{±$.103$} \\
\midrule
\multirow{4}{*}{Combinations} & Simple + Multiple & $0.29$ \tiny{±$.071$} & $0.36$ \tiny{±$.073$} & $0.40$ \tiny{±$.075$} & \underline{$0.42$} \tiny{±$.074$} & $0.38$ \tiny{±$.076$} & \best{$0.44$} \tiny{±$.075$} & $0.38$ \tiny{±$.075$} & $0.30$ \tiny{±$.072$} & $0.26$ \tiny{±$.073$} & $0.07$ \tiny{±$.066$} \\
 & Simple + Parallel & $0.25$ \tiny{±$.069$} & $0.29$ \tiny{±$.070$} & $0.37$ \tiny{±$.072$} & \underline{$0.40$} \tiny{±$.070$} & $0.35$ \tiny{±$.072$} & \best{$0.41$} \tiny{±$.072$} & $0.35$ \tiny{±$.071$} & $0.24$ \tiny{±$.069$} & $0.23$ \tiny{±$.068$} & $0.09$ \tiny{±$.061$} \\
 & Multiple + Parallel-Multiple & $0.29$ \tiny{±$.076$} & $0.27$ \tiny{±$.078$} & $0.41$ \tiny{±$.075$} & \best{$0.46$} \tiny{±$.076$} & $0.39$ \tiny{±$.079$} & \underline{$0.45$} \tiny{±$.077$} & $0.39$ \tiny{±$.079$} & $0.24$ \tiny{±$.080$} & $0.28$ \tiny{±$.085$} & $0.17$ \tiny{±$.073$} \\
 & All Combined & $0.29$ \tiny{±$.050$} & $0.29$ \tiny{±$.051$} & $0.40$ \tiny{±$.051$} & \best{$0.44$} \tiny{±$.051$} & $0.38$ \tiny{±$.052$} & \underline{$0.43$} \tiny{±$.052$} & $0.38$ \tiny{±$.052$} & $0.25$ \tiny{±$.052$} & $0.24$ \tiny{±$.053$} & $0.15$ \tiny{±$.048$} \\
\midrule
\multirow{2}{*}{Unanswerable} & Simple + Irrelevance & $0.50$ \tiny{±$.055$} & $0.58$ \tiny{±$.049$} & \underline{$0.59$} \tiny{±$.051$} & \best{$0.61$} \tiny{±$.050$} & $0.54$ \tiny{±$.053$} & $0.52$ \tiny{±$.055$} & $0.55$ \tiny{±$.052$} & $0.55$ \tiny{±$.053$} & $0.44$ \tiny{±$.055$} & $0.03$ \tiny{±$.055$} \\
 & All Combined + Irrelevance & $0.39$ \tiny{±$.040$} & $0.46$ \tiny{±$.039$} & \underline{$0.48$} \tiny{±$.039$} & \best{$0.52$} \tiny{±$.039$} & $0.45$ \tiny{±$.041$} & $0.45$ \tiny{±$.041$} & $0.45$ \tiny{±$.040$} & $0.43$ \tiny{±$.042$} & $0.32$ \tiny{±$.043$} & $0.01$ \tiny{±$.036$} \\
\bottomrule
\end{tabular}
}
\caption{Results: Mean PRR@30\% $\pm$ mean standard error (mean over eight LLMs) for individual tasks, task combinations, and unanswerable requests. PRR@30\% is the area-based Prediction Rejection Ratio integrated up to 30\% rejection (oracle $=1$, random $=0$; higher is better). Best result in \best{green}, \underline{underlined} second best.}
\label{tab:prr30}
\end{table}

\begin{table}[h!]
\centering
\scalebox{0.5}{
\begin{tabular}{llcccccccccc}
\toprule
& & \multicolumn{4}{c}{\textbf{Single-sample}} & \multicolumn{4}{c}{\textbf{Multi-sample}} & \multicolumn{2}{c}{\textbf{Other}} \\
\cmidrule(lr){3-6} \cmidrule(lr){7-10} \cmidrule(lr){11-12}
\textbf{Category} & \textbf{Task} & \textbf{MAX} & \textbf{AVG} & \textbf{G-NLL} & \textbf{G-NLL}$_{\mathrm{SMT}}$ & $\mathbf{SE}_{EXM}$ & $\mathbf{SE}_{AST}$ & $\mathbf{DSE}_{EXM}$ & \textbf{PE} & $\mathbf{P}_{true}$ & \textbf{LEN} \\
\midrule
\multirow{4}{*}{Individual} & Simple & $0.41$ \tiny{±$.092$} & $0.47$ \tiny{±$.092$} & \underline{$0.49$} \tiny{±$.092$} & \best{$0.53$} \tiny{±$.091$} & $0.44$ \tiny{±$.103$} & \underline{$0.49$} \tiny{±$.101$} & $0.43$ \tiny{±$.103$} & $0.38$ \tiny{±$.095$} & $0.35$ \tiny{±$.100$} & $0.07$ \tiny{±$.094$} \\
 & Multiple & $0.36$ \tiny{±$.132$} & $0.43$ \tiny{±$.135$} & \underline{$0.46$} \tiny{±$.137$} & \best{$0.47$} \tiny{±$.135$} & $0.41$ \tiny{±$.145$} & $0.45$ \tiny{±$.143$} & $0.41$ \tiny{±$.147$} & $0.37$ \tiny{±$.139$} & $0.30$ \tiny{±$.142$} & $0.16$ \tiny{±$.141$} \\
 & Parallel & $0.26$ \tiny{±$.115$} & \underline{$0.38$} \tiny{±$.113$} & $0.37$ \tiny{±$.115$} & \best{$0.41$} \tiny{±$.114$} & $0.30$ \tiny{±$.120$} & $0.35$ \tiny{±$.124$} & $0.30$ \tiny{±$.120$} & $0.29$ \tiny{±$.120$} & $0.21$ \tiny{±$.118$} & $0.01$ \tiny{±$.106$} \\
 & Parallel-Multiple & $0.37$ \tiny{±$.104$} & $0.44$ \tiny{±$.104$} & \underline{$0.46$} \tiny{±$.105$} & \best{$0.52$} \tiny{±$.100$} & $0.41$ \tiny{±$.110$} & \underline{$0.46$} \tiny{±$.106$} & $0.40$ \tiny{±$.112$} & $0.37$ \tiny{±$.113$} & $0.31$ \tiny{±$.119$} & $0.08$ \tiny{±$.111$} \\
\midrule
\multirow{4}{*}{Combinations} & Simple + Multiple & $0.39$ \tiny{±$.074$} & $0.45$ \tiny{±$.074$} & $0.48$ \tiny{±$.075$} & \best{$0.51$} \tiny{±$.075$} & $0.43$ \tiny{±$.082$} & \underline{$0.49$} \tiny{±$.081$} & $0.42$ \tiny{±$.082$} & $0.38$ \tiny{±$.077$} & $0.33$ \tiny{±$.080$} & $0.10$ \tiny{±$.079$} \\
 & Simple + Parallel & $0.34$ \tiny{±$.073$} & $0.36$ \tiny{±$.074$} & \underline{$0.43$} \tiny{±$.073$} & \best{$0.47$} \tiny{±$.071$} & $0.37$ \tiny{±$.078$} & $0.42$ \tiny{±$.078$} & $0.37$ \tiny{±$.077$} & $0.29$ \tiny{±$.076$} & $0.30$ \tiny{±$.073$} & $0.13$ \tiny{±$.073$} \\
 & Multiple + Parallel-Multiple & $0.38$ \tiny{±$.078$} & $0.35$ \tiny{±$.082$} & \underline{$0.48$} \tiny{±$.079$} & \best{$0.53$} \tiny{±$.075$} & $0.44$ \tiny{±$.084$} & \underline{$0.48$} \tiny{±$.082$} & $0.43$ \tiny{±$.084$} & $0.30$ \tiny{±$.084$} & $0.33$ \tiny{±$.088$} & $0.23$ \tiny{±$.082$} \\
 & All Combined & $0.37$ \tiny{±$.052$} & $0.36$ \tiny{±$.053$} & \underline{$0.46$} \tiny{±$.052$} & \best{$0.51$} \tiny{±$.051$} & $0.41$ \tiny{±$.055$} & $0.45$ \tiny{±$.056$} & $0.41$ \tiny{±$.056$} & $0.30$ \tiny{±$.055$} & $0.31$ \tiny{±$.057$} & $0.19$ \tiny{±$.054$} \\
\midrule
\multirow{2}{*}{Unanswerable} & Simple + Irrelevance & $0.50$ \tiny{±$.054$} & $0.56$ \tiny{±$.051$} & \underline{$0.57$} \tiny{±$.052$} & \best{$0.60$} \tiny{±$.049$} & $0.50$ \tiny{±$.054$} & $0.47$ \tiny{±$.056$} & $0.52$ \tiny{±$.054$} & $0.53$ \tiny{±$.053$} & $0.47$ \tiny{±$.052$} & $0.02$ \tiny{±$.056$} \\
 & All Combined + Irrelevance & $0.42$ \tiny{±$.041$} & $0.48$ \tiny{±$.041$} & \underline{$0.50$} \tiny{±$.040$} & \best{$0.55$} \tiny{±$.039$} & $0.45$ \tiny{±$.043$} & $0.44$ \tiny{±$.043$} & $0.45$ \tiny{±$.042$} & $0.45$ \tiny{±$.042$} & $0.37$ \tiny{±$.043$} & $-0.01$ \tiny{±$.040$} \\
\bottomrule
\end{tabular}
}
\caption{Results: Mean PRR@50\% $\pm$ mean standard error (mean over eight LLMs) for individual tasks, task combinations, and unanswerable requests. PRR@50\% is the area-based Prediction Rejection Ratio integrated up to 50\% rejection (oracle $=1$, random $=0$; higher is better). Best result in \best{green}, \underline{underlined} second best.}
\label{tab:prr50}
\end{table}


